\section{第三题综合模拟专题}
\subsection{长题意的实现流程}
先写一张规则表：输入是什么、状态存在哪里、操作改变什么、查询输出什么、异常如何处理。再按职责拆成解析层、状态层和执行层。只有确认每层的输入输出后，再连成完整程序。下面五个案例规则完整，分别训练路径、映射状态、递归解析、树形结构和依赖展开，可以脱离外部题面学习。

\example{路径规范化}
\textbf{题目。}给定固定的当前目录cwd和q条路径。以斜杠开头的是绝对路径，否则相对于cwd。连续斜杠视为一个；点号目录表示当前目录；两个点号表示父目录，在根目录继续上移仍为根。普通目录名只含字母、数字、下划线。路径不含空格且非空。

输出规范绝对路径：以一个斜杠开头，不含点号、两个点号或重复斜杠；除根目录外不以斜杠结尾。\textbf{每条查询独立使用原来的cwd，不改变当前目录。}路径长度不超过 $10^4$，全部输入路径总长度不超过 $2\times10^5$，$1\le q\le10^4$。第一行cwd，第二行q，随后q行各一条路径。cwd保证为绝对路径，也允许特殊片段。
\begin{verbatim}
输入
/a/b
5
../c
/x//./y/../
../../../
.
tmp/../../z
输出
/a/c
/x
/
/a/b
/a/z
\end{verbatim}
\textbf{解答。}用vector保存从根向下的目录名。普通片段入栈，两个点号在非空时弹栈，空片段和点号忽略。绝对路径从空栈开始，相对路径从cwd的副本开始。

最后一条查询从 $[a,b]$ 出发，读tmp得到 $[a,b,tmp]$，两次父目录操作后为 $[a]$，加入z后得到 $[a,z]$。不变量是：栈中的目录始终表示已经处理的路径前缀的规范形式。

设cwd长度为B，第i个查询长度为 $L_i$。复制基目录的总成本最坏为 $O(qB)$，加上扫描与输出为 $O(qB+\sum L_i)$。不能漏算复制和输出中重复的cwd长度。
\Code{19_path.cpp}
\textbf{调试。}先测根目录与普通绝对路径，再测相对路径和多次上移，最后检查是否错误地修改了共享cwd。若后续查询受前面查询影响，优先检查是否复用了同一可变目录数组。

\example{授权和撤销的状态管理}
\textbf{题目。}执行q条命令：\texttt{GRANT user resource}授予权限；\texttt{REVOKE user resource}撤销；\texttt{CHECK user resource}查询并输出YES或NO。重复授予不叠加；撤销不存在的权限不报错；未出现用户默认无权限。操作依次执行，只有CHECK输出。$1\le q\le2\times10^5$，名称为1到20个字母或数字。本题不含角色继承、通配符或优先级。
\begin{verbatim}
输入
8
CHECK alice file1
GRANT alice file1
GRANT alice file1
CHECK alice file1
REVOKE alice file1
CHECK alice file1
REVOKE bob file2
CHECK bob file2
输出
NO
YES
NO
NO
\end{verbatim}
\textbf{解答。}外层map按用户名定位，内层set保存资源。set处理重复授权。查询和撤销用find，避免无意创建空记录。所有操作依次执行，因此CHECK反映此前的状态。

设用户数U，每位用户最多持有R项权限，单次比较次数为 $O(\log U+\log R)$，还需考虑字符串比较长度。总存储不超过授予过的记录规模。权限是用户与资源的二元关系，不能只判断用户是否出现过；撤销file1不能影响file2。
\Code{20_permissions.cpp}

\example{带括号的整数表达式}
\textbf{题目。}输入一行表达式，支持非负整数字面量、加、减、乘、括号和一元正负号。乘法优先于加减，同级二元运算从左到右；一元符号优先于乘法；只忽略普通空格。合法时输出结果，语法错误输出ERROR。不支持除法、变量和隐式乘法，例如 \texttt{2(3+4)} 非法。

行长不超过1000，括号与连续一元符号嵌套不超过100层，字面量不超过 $10^6$；保证求值过程中涉及的数值运算的中间结果绝对值不超过 $10^{12}$，本题只讨论语法错误，不要求处理溢出。
\begin{verbatim}
输入
2 * (3 + 4) - -5
输出
19
另一组输入
1 + (2 * 3
另一组输出
ERROR
\end{verbatim}
\textbf{按语法拆函数。}
\begin{verbatim}
expression = term { ('+' | '-') term }
term       = factor { '*' factor }
factor     = number | '(' expression ')' | '+' factor | '-' factor
\end{verbatim}
expression负责加减，term负责乘法，factor负责括号、数字与一元符号。共享指针p指向下一个未处理字符，每个函数返回其负责部分的值。term先完整处理乘法，从而保证优先级；循环保证减法左结合，例如 $8-3-2=3$。

解析前面一个合法表达式不代表整行合法。例如 \texttt{1 2}，解析1后还剩2，必须输出ERROR。最后检查是否读到行尾。时间 $O(L)$，递归空间与嵌套深度有关。
\Code{21_expression.cpp}
\textbf{调试。}依次验证单数字、一元负号、连续加减、乘法优先级、嵌套括号，再测空串、缺右括号、多右括号和多余字符。查看每个函数消费到哪里，比只看最终错误数值更有帮助。

\example{树形文档的后代选择查询}
\textbf{题目。}给出按先序排列的一片有根森林，共n个节点。每行记录深度、标签和标识。根深度0，孩子比父亲深1；输入结构合法，每次深度最多增加1。标识为减号表示无标识，其余标识互不相同；标签和标识只含字母数字，长度1到20。

查询由空格分隔的选择符组成。普通选择符匹配标签，以井号开头的选择符匹配标识。\texttt{div p}表示匹配p节点，且它的某个\textbf{严格祖先}是div；\texttt{div div p}需要两个不同的div祖先按顺序出现，允许隔着其他节点。节点不能把自己当祖先。

$1\le n\le500$，$1\le q\le50$，最大深度100，每条查询1到10个选择符，查询行总长度不超过1000。第一行n、q，随后n行节点，最后q行查询。每次输出匹配数和按输入顺序排列的节点行号，编号从1开始。
\begin{verbatim}
输入
7 5
0 html -
1 body -
2 div main
3 p one
2 div side
3 p two
2 p three
p
div p
#main p
div div p
#side
输出
3 4 6 7
2 4 6
1 4
0
1 5
\end{verbatim}
\textbf{建树。}用祖先栈。读取深度d的节点时，把栈截断到长度d；栈顶就是父亲，空栈则为根，然后压入自己。读取同深度兄弟时，旧兄弟会被弹出。

\textbf{匹配。}从查询末尾开始，先检查候选节点本身，再沿父链逆序匹配剩余选择符。找到一个就继续寻找更高祖先。选择最近的匹配祖先不会损失可行解，因为它给剩余前缀保留最多的上方祖先。找齐全部选择符则接受。

候选检查最多走树高H条父边并处理K个选择符，整体为 $O(qn(H+K))$，本题规模下足够。节点7虽是p，但没有div祖先；若程序把它算入 \texttt{div p}，通常是把“之前出现的节点”误认为“当前祖先”。
\Code{22_selector.cpp}

\example{模板展开与循环依赖}
\textbf{题目。}有n个命名模板，每个由若干以空白分隔的词组成。以at符号开头的词引用另一个模板，其余词直接输出。定义全部读完后，展开指定起始模板，以单个空格连接所有普通词。从起始模板可达的引用若出现未定义名称或环，输出ERROR；不可达模板的错误不影响本次输出。

模板名互不重复，只含字母数字，长度不超过20；引用符号后必须有名称。允许空模板。$1\le n\le100$，每个定义最多100个词，每词长度最多20。保证无环情况下每个模板完整展开不超过2000个词。第一行n，随后每个定义包含名称、词数k及k个词，最后给出起始名称。
\begin{verbatim}
输入
3
title 1 world
line 3 hello @title !
all 2 @line @line
all
输出
hello world ! hello world !
另一组输入
2
a 1 @b
b 1 @a
a
另一组输出
ERROR
\end{verbatim}
\textbf{三色状态。}0表示未算，1表示正在递归，2表示已经完成并缓存。访问状态1说明当前依赖链回到了自己；访问状态2直接复用。第二次引用同一模板不必然成环，关键看第一次展开是否结束。

\textbf{解答。}先读入全部定义以允许前向引用。按词序展开，普通词加入结果，引用词接上被引用模板的结果。记忆化避免重算，但每次拼接仍复制相应词。时间取决于查找及所有实际复制的词数；本题每个模板最多缓存2000词，总缓存约不超过 $100\times2000$ 词，不能忽略输出规模而声称总是O(n)。
\Code{23_templates.cpp}

\exercise
\textbf{练习1。}cwd为根，路径为 \texttt{../../a}，结果是什么？\textbf{解答。}\texttt{/a}，超出根的上移被忽略。

\textbf{练习2。}alice有A和B，撤销A后查询B输出什么？\textbf{解答。}YES，资源权限分别维护。

\textbf{练习3。}\texttt{-2*-3+4}的值是多少？\textbf{解答。}10。一元负号属于factor，乘法先得6，再加4。

\textbf{练习4。}一个div节点能否独自匹配 \texttt{div div}？\textbf{解答。}不能，两个选择符必须匹配不同节点，前一个位于严格祖先链。

\textbf{练习5。}A引用B两次，B只含x，是否有环？\textbf{解答。}没有。第一次结束后B为状态2，第二次复用，输出 \texttt{x x}。
\newpage
